% GENERATED FILE. Edit canonical lessons, then run npm run generate:books.
% canonical-source-hash: fnv1a64:9cdacddbd3a2fe74
% canonical-lessons: MR-C426-rokad, MR-C426-dhanades, MR-C426-cillar, MR-C426-rakkam, MR-C426-jama-karne

\chapter{Cash and Cheques}
\label{ch:mr-cash-and-cheques}
\hlchaptermodality{426}

\begin{chapteropening}
\textbf{By the end of this chapter:} \emph{I can talk about cash, cheques, small change and sums of money, and paying money in.}

Complete the last lesson of chapter 426: I can talk about cash, cheques, small change and sums of money, and paying money in.
\end{chapteropening}

\section[\texorpdfstring{\mr{रोकड} (rokaḍ)}{रोकड (rokaḍ)}]{\mr{रोकड} (rokaḍ) — cash}

\label{lesson:MR-C426-rokad}



\begin{quote}
Before the new one: say the Marathi for fake, then the Marathi for a copy.
\end{quote}

\subsection*{You'll want to know: \mr{रोकड}}
\textbf{\mr{रोकड}} — \emph{rokaḍ} — \textquotedblleft{}cash\textquotedblright{}.

\begin{quote}
\mr{दुकानदाराला} \mr{रोकड} \mr{हवी} \mr{होती}.
\end{quote}

\begin{quote}
\emph{dukāndārālā rokaḍ havī hotī.}
\end{quote}

\begin{quote}
The shopkeeper wanted cash.
\end{quote}

Feminine.

\mr{रोख} is the adjective: \mr{रोख} \mr{पैसे}, cash money.

\begin{cousinweb}
\textbf{\mr{रोकड}} (\emph{rokaḍ}) is inherited: CDIAL lists Marathi \emph{rokaḍ}, \textquotedblleft{}cash\textquotedblright{}, under \emph{*rōkka}, \textquotedblleft{}ready money\textquotedblright{} (CDIAL 10828).
\end{cousinweb}

\begin{grammarlens}[title={What you've built}]
One more word for getting around.
\end{grammarlens}

\subsection*{Your turn}
\begin{itemize}
\raggedright
  \item \emph{Say it:} \emph{rokaḍ}
  \item \emph{Say it:} \emph{rokaḍ}, once more
  \item \emph{Say it:} \emph{dukāndārālā rokaḍ havī hotī.} — The shopkeeper wanted cash.
  \item \emph{Recall:} say the Marathi for fake, then the Marathi for a copy, then say \emph{rokaḍ} again
\end{itemize}

\begin{tcolorbox}[breakable,colback=teal!4,colframe=teal!35!black,title={Before you move on}]
What does \mr{रोकड} mean? (\textquotedblleft{}Cash\textquotedblright{}.) Say it once more.
\end{tcolorbox}
\section[\texorpdfstring{\mr{धनादेश} (dhanādeś)}{धनादेश (dhanādeś)}]{\mr{धनादेश} (dhanādeś) — a cheque}

\label{lesson:MR-C426-dhanades}



\begin{quote}
Before the new one: say the Marathi for a copy, then the Marathi for cash.
\end{quote}

\subsection*{You'll want to know: \mr{धनादेश}}
\textbf{\mr{धनादेश}} — \emph{dhanādeś} — \textquotedblleft{}a cheque\textquotedblright{}.

\begin{quote}
\mr{त्याने} \mr{भाड्यासाठी} \mr{धनादेश} \mr{दिला}.
\end{quote}

\begin{quote}
\emph{tyāne bhāḍyāsāṭhī dhanādeś dilā.}
\end{quote}

\begin{quote}
He gave a cheque for the rent.
\end{quote}

Masculine.

A formal word; in speech many people say \mr{चेक}.

\begin{cousinweb}
\textbf{\mr{धनादेश}} (\emph{dhanādeś}) is a modern word built from Sanskrit \emph{dhana}, \textquotedblleft{}wealth\textquotedblright{}, and \emph{ādeśa}, \textquotedblleft{}an order\textquotedblright{} (Monier-Williams).
\end{cousinweb}

\begin{grammarlens}[title={What you've built}]
One more word for getting around.
\end{grammarlens}

\subsection*{Your turn}
\begin{itemize}
\raggedright
  \item \emph{Say it:} \emph{dhanādeś}
  \item \emph{Say it:} \emph{dhanādeś}, once more
  \item \emph{Say it:} \emph{tyāne bhāḍyāsāṭhī dhanādeś dilā.} — He gave a cheque for the rent.
  \item \emph{Recall:} say the Marathi for a copy, then the Marathi for cash, then say \emph{dhanādeś} again
\end{itemize}

\begin{tcolorbox}[breakable,colback=teal!4,colframe=teal!35!black,title={Before you move on}]
What does \mr{धनादेश} mean? (\textquotedblleft{}A cheque\textquotedblright{}.) Say it once more.
\end{tcolorbox}
\section[\texorpdfstring{\mr{चिल्लर} (cillar)}{चिल्लर (cillar)}]{\mr{चिल्लर} (cillar) — small change}

\label{lesson:MR-C426-cillar}



\begin{quote}
Before the new one: say the Marathi for cash, then the Marathi for a cheque.
\end{quote}

\subsection*{You'll want to know: \mr{चिल्लर}}
\textbf{\mr{चिल्लर}} — \emph{cillar} — \textquotedblleft{}small change\textquotedblright{}.

\begin{quote}
\mr{माझ्याकडे} \mr{चिल्लर} \mr{नाही}.
\end{quote}

\begin{quote}
\emph{mājhyākaḍe cillar nāhī.}
\end{quote}

\begin{quote}
I have no small change.
\end{quote}

Feminine.

Coins and small notes; it also means petty, trivial.

\begin{cousinweb}
\textbf{\mr{चिल्लर}} (\emph{cillar}) is a Marathi word for loose coins; its origin is not settled.
\end{cousinweb}

\begin{grammarlens}[title={What you've built}]
One more word for getting around.
\end{grammarlens}

\subsection*{Your turn}
\begin{itemize}
\raggedright
  \item \emph{Say it:} \emph{cillar}
  \item \emph{Say it:} \emph{cillar}, once more
  \item \emph{Say it:} \emph{mājhyākaḍe cillar nāhī.} — I have no small change.
  \item \emph{Recall:} say the Marathi for cash, then the Marathi for a cheque, then say \emph{cillar} again
\end{itemize}

\begin{tcolorbox}[breakable,colback=teal!4,colframe=teal!35!black,title={Before you move on}]
What does \mr{चिल्लर} mean? (\textquotedblleft{}Small change\textquotedblright{}.) Say it once more.
\end{tcolorbox}
\section[\texorpdfstring{\mr{रक्कम} (rakkam)}{रक्कम (rakkam)}]{\mr{रक्कम} (rakkam) — a sum of money}

\label{lesson:MR-C426-rakkam}



\begin{quote}
Before the new one: say the Marathi for a cheque, then the Marathi for small change.
\end{quote}

\subsection*{You'll want to know: \mr{रक्कम}}
\textbf{\mr{रक्कम}} — \emph{rakkam} — \textquotedblleft{}a sum of money\textquotedblright{}.

\begin{quote}
\mr{ही} \mr{रक्कम} \mr{खूप} \mr{मोठी} \mr{आहे}.
\end{quote}

\begin{quote}
\emph{hī rakkam khūp moṭhī āhe.}
\end{quote}

\begin{quote}
This sum is very large.
\end{quote}

Feminine.

A feminine noun: \mr{मोठी} \mr{रक्कम}.

\begin{cousinweb}
\textbf{\mr{रक्कम}} (\emph{rakkam}) is a loan through Persian, usually traced to Arabic \emph{raqam}, \textquotedblleft{}a written figure\textquotedblright{}.
\end{cousinweb}

\begin{grammarlens}[title={What you've built}]
One more word for getting around.
\end{grammarlens}

\subsection*{Your turn}
\begin{itemize}
\raggedright
  \item \emph{Say it:} \emph{rakkam}
  \item \emph{Say it:} \emph{rakkam}, once more
  \item \emph{Say it:} \emph{hī rakkam khūp moṭhī āhe.} — This sum is very large.
  \item \emph{Recall:} say the Marathi for a cheque, then the Marathi for small change, then say \emph{rakkam} again
\end{itemize}

\begin{tcolorbox}[breakable,colback=teal!4,colframe=teal!35!black,title={Before you move on}]
What does \mr{रक्कम} mean? (\textquotedblleft{}A sum of money\textquotedblright{}.) Say it once more.
\end{tcolorbox}
\section[\texorpdfstring{\mr{जमा} \mr{करणे} (jamā karṇe)}{जमा करणे (jamā karṇe)}]{\mr{जमा} \mr{करणे} (jamā karṇe) — to deposit; to collect}

\label{lesson:MR-C426-jama-karne}



\begin{quote}
Before the new one: say the Marathi for small change, then the Marathi for a sum of money.
\end{quote}

\subsection*{You'll want to know: \mr{जमा} \mr{करणे}}
\textbf{\mr{जमा} \mr{करणे}} — \emph{jamā karṇe} — \textquotedblleft{}to deposit; to collect\textquotedblright{}.

\begin{quote}
\mr{मी} \mr{पगाराचे} \mr{पैसे} \mr{खात्यात} \mr{जमा} \mr{केले}.
\end{quote}

\begin{quote}
\emph{mī pagārāce paise khātyāt jamā kele.}
\end{quote}

\begin{quote}
I paid my salary into my account.
\end{quote}

Past, with -\mr{ने} (he …): \mr{त्याने} \mr{जमा} \mr{केलं}, \emph{tyāne jamā kelaṁ}.

The verb agrees with the money: \mr{पैसे} \mr{जमा} \mr{केले}, \mr{रक्कम} \mr{जमा} \mr{केली}.

\begin{cousinweb}
\textbf{\mr{जमा} \mr{करणे}} (\emph{jamā karṇe}) joins \textbf{\mr{जमा}}, \textquotedblleft{}collected, deposited\textquotedblright{}, a loan through Persian, usually traced to Arabic \emph{jamʿ}, \textquotedblleft{}gathering\textquotedblright{}, to the light verb \textbf{\mr{करणे}}.
\end{cousinweb}

\begin{grammarlens}[title={What you've built}]
One more word for getting around.
\end{grammarlens}

\subsection*{Your turn}
\begin{itemize}
\raggedright
  \item \emph{Say it:} \emph{jamā karṇe}
  \item \emph{Say it:} \emph{jamā karṇe}, once more
  \item \emph{Say it:} \emph{mī pagārāce paise khātyāt jamā kele.} — I paid my salary into my account.
  \item \emph{Recall:} say the Marathi for small change, then the Marathi for a sum of money, then say \emph{jamā karṇe} again
\end{itemize}

\begin{tcolorbox}[breakable,colback=teal!4,colframe=teal!35!black,title={Before you move on}]
What does \mr{जमा} \mr{करणे} mean? (\textquotedblleft{}To deposit; to collect\textquotedblright{}.) Say it once more.
\end{tcolorbox}
